\documentclass[10pt, a4paper]{article}

% ---- Универсальный одностраничный шаблон резюме. Замените ВЫМЫШЛЕННУЮ персону
% ---- "Alex Doe" и все заглушки на свои данные. Строки с "% ATS-*" переписывает
% ---- resume/tailor.js: оставьте их там, где нужен вставляемый контент.

\usepackage[
    ignoreheadfoot,
    top=1.8cm, bottom=1.8cm, left=1.8cm, right=1.8cm,
    footskip=0.8cm,
]{geometry}
\usepackage{titlesec}
\usepackage{array}
\usepackage[dvipsnames]{xcolor}
\definecolor{primaryColor}{RGB}{0, 0, 0}
\usepackage{enumitem}
\usepackage{amsmath}
\usepackage[
    pdftitle={Alex Doe CV},
    pdfauthor={Alex Doe},
    colorlinks=true,
    urlcolor=primaryColor
]{hyperref}
\usepackage{calc}
\usepackage{paracol}
\usepackage{iftex}

\ifPDFTeX
    \input{glyphtounicode}
    \pdfgentounicode=1
    \usepackage[T1,T2A]{fontenc}
    \usepackage[utf8]{inputenc}
    \usepackage[english,russian]{babel}
    \usepackage{lmodern}
\fi

\renewcommand{\familydefault}{\ttdefault}
\raggedright
\pagestyle{empty}
\setcounter{secnumdepth}{0}
\setlength{\parindent}{0pt}
\setlength{\columnsep}{0.15cm}
\pagenumbering{gobble}

\titleformat{\section}{\bfseries\large}{}{0pt}{}[\vspace{0.5pt}\titlerule]
\titlespacing{\section}{-1pt}{0.25cm}{0.15cm}

\newenvironment{highlights}{
    \begin{itemize}[topsep=0.08cm, parsep=0.08cm, partopsep=0pt, itemsep=0pt, leftmargin=10pt]
}{
    \end{itemize}
}
\newenvironment{onecolentry}{\begin{adjustwidth}{0.00001cm}{0.00001cm}}{\end{adjustwidth}}
\newenvironment{twocolentry}[2][]{
    \onecolentry
    \def\secondColumn{#2}
    \setcolumnwidth{\fill, 4.5cm}
    \begin{paracol}{2}
}{
    \switchcolumn \raggedleft \secondColumn
    \end{paracol}
    \endonecolentry
}
\usepackage{changepage}
\newenvironment{header}{\setlength{\topsep}{0pt}\par\kern\topsep\centering\linespread{1.3}}{\par\kern\topsep}

\begin{document}
    \newcommand{\AND}{\unskip\cleaders\copy\ANDbox\hskip\wd\ANDbox\ignorespaces}
    \newsavebox\ANDbox
    \sbox\ANDbox{$|$}

    \begin{header}
        \fontsize{25pt}{25pt}\selectfont Alex Doe

        \vspace{5pt}
        \normalsize
        \mbox{\href{https://example.com/alex-doe}{example.com/alex-doe}}%
        \kern 5.0pt \AND \kern 5.0pt%
        \mbox{\href{mailto:alex.doe@example.com}{alex.doe@example.com}}%
        \kern 5.0pt \AND \kern 5.0pt%
        \mbox{\href{https://www.linkedin.com/in/your-handle}{linkedin.com/in/your-handle}}%
    \end{header}

    \vspace{5pt - 0.3cm}

    % ATS-HEADLINE (tailor.js подставит следующую строку, если передан headline)
    \section{Обо мне}

    \begin{onecolentry}
        TODO: 2-3 предложения своими словами. Кто вы, что делаете, за какой
        результат вас знают. Только то, что сможете защитить на интервью.
    \end{onecolentry}

    \vspace{0.15cm}
    \section{Опыт}

    \begin{twocolentry}{Мес 20XX -- н.в.}
        \textbf{Ваша последняя должность}, Компания -- Удалённо\end{twocolentry}
    \vspace{0.08cm}
    \begin{onecolentry}
        \begin{highlights}
            \item TODO: результат, за который отвечали, с метрикой.
            \item TODO: второй пункт: масштаб, решение, эффект.
            \item TODO: третий пункт, ближе к целевой роли.
        \end{highlights}
    \end{onecolentry}

    \vspace{0.15cm}
    \begin{twocolentry}{Мес 20XX -- Мес 20XX}
        \textbf{Предыдущая должность}, Компания -- Удалённо\end{twocolentry}
    \vspace{0.08cm}
    \begin{onecolentry}
        \begin{highlights}
            \item TODO: результат.
            \item TODO: результат.
        \end{highlights}
    \end{onecolentry}

    \vspace{0.15cm}
    \section{Образование}

    \begin{twocolentry}{Мес 20XX -- Мес 20XX}
        \textbf{Степень}, Направление
    \end{twocolentry}
    \vspace{0.05cm}
    \begin{onecolentry}
        Учебное заведение
    \end{onecolentry}

    \vspace{0.15cm}
    \section{Навыки}

    \begin{onecolentry}
        % ATS-SKILLS-LINE (tailor.js подставит следующую строку из матча навыков)
        \textbf{Product:} Strategy \textperiodcentered{} Discovery \textperiodcentered{} Roadmap \textbf{|} \textbf{Tools:} Jira \textperiodcentered{} Notion
    \end{onecolentry}

\end{document}
